\chapter{Discussion}
% ===================================================================

\section{Relation to Existing Work}

The logical topology on the bisimulation quotient is standard in
domain theory and coalgebra~\cite{abramsky1991domain}.
Its use here as a substrate for \emph{epistemic} topology ---
the topology of what an agent can know --- is less standard,
though related to the topological approach to common knowledge
in~\cite{fagin1995reasoning}.

The compactness-as-consensus argument is new, to the authors'
knowledge.
It is inspired by the compactness theorem in logic (if every finite
subset of a theory has a model, the whole theory does) but runs in
the opposite direction.

The nerve construction is standard in algebraic
topology~\cite{hatcher2002algebraic}.
Its use as a model of inter-community ontological structure appears
to be new in this computational setting.

The connection to effective field theory and the renormalization group
is conjectural but suggestive.
The renormalization group as a morphism in a category of field
theories has been studied in the physics
literature~\cite{wilson1974renormalization};
whether it can be modeled as a morphism in $\mathbf{GSLT}$ is
an open question.

\section{The Deeper Implication}

The argument of this paper, if completed, would establish something
philosophically significant: that the hierarchy of physical scales ---
quarks, hadrons, nuclei, atoms, molecules, and so on --- is not a
contingent feature of our particular universe.
It is the \emph{inevitable structure} of any sufficiently large
population of interacting processes, arising from the communication
geometry of the GSLT and the compactness condition for consensus.

Any civilization of scientists, anywhere in the computational
universe, will encounter the same kind of hierarchical structure
in their experiments.
They will build correct theories of it at each scale.
What they will not discover --- without resources that grow with
the scale of the population --- is the underlying bisimulation
quotient from which the hierarchy emerges.

This is not a counsel of despair.
The scientific method is not broken; it is working exactly as
it should.
The cluster hierarchy is real, and the theories that describe it
are true.
The deeper ontology --- the bisimulation quotient with its prime
combinators --- is not a more correct description that the
scientists are failing to reach.
It is a description at a different resolution, requiring
experimental resources that scale with the size of the full
population.

The implication for artificial intelligence is direct.
An AI system, however powerful, is itself a term in some GSLT
and therefore subject to the same budget constraints and
communication geometry as any other agent.
It will build theories of the cluster hierarchy it inhabits.
Whether it can transcend that hierarchy --- whether it can, with
sufficient resources, push the scientific method toward the
bisimulation quotient --- is a question about the relationship
between its budget and the scale of the population it is embedded in.
That question is open.

% ===================================================================
